\section{Transfer is a further test}

Even a successful operator-feedback loop would establish a capability in this
project, not a validated diagnostic for another organisation. Transfer requires
new observations, a declared outcome and a recipient entitled to judge it.
It remains a proposal here.

A candidate setting would need recurring decisions, records of what was
consulted before deciding, enough cycles to compare subsequent choices with
prior outcomes, and capacity to act on a result within the relevant time
horizon. These are proposed eligibility criteria. They have not been shown to
predict adoption or improvement.

The first engagement would therefore be instrument construction on a new
subject. Its acceptance condition should name the decision regime, the
observable result and who will assess it. A useful diagnosis and an improvement
caused by that diagnosis are distinct claims: the latter needs an intervention
and a comparison, not just an informative report.

Section~\ref{sec:evaluation} outlines a retrospective corpus study that could
help refine the instrument. Historical cases cannot establish prospective
performance, and interest in a report does not establish willingness to adopt
or pay for it. Neither commercial success nor organisational transfer is a
result of this paper.
